%% =============================================================================
%% sm_c.tex — Supplementary Material C: Survey Design Theory
%% Body content only (subsections); \section and \label defined in osm_body.tex
%% Primary source: dev/log/06_research-notes-survey-design-pseudo-posterior.qmd
%%   Parts 2--4 (pseudo-posterior theory, sandwich, Cholesky, DER)
%% Corrections applied: T1-1 (BvM reframe), T1-4 (MCMC vs ESS remark),
%%   T2-1 (block-wise Cholesky caveat), T2-2 (DER heuristic),
%%   T2-3 ("exactly" -> "in expectation"), T2-4 (intercept identifiability)
%% =============================================================================

This appendix develops the survey design theory underlying the sandwich
variance estimator and the Cholesky affine transformation described in
\cref{sec:model}. We begin with the formal pseudo-likelihood framework
and the proof that the pseudo-posterior is proper
(\cref{thm:propriety}), establish the Bernstein--von Mises conditions
and their finite-sample diagnostic for the NSECE weights, derive the
cluster-robust sandwich estimator and its block structure, provide the
full proof of the Cholesky affine transformation
(\cref{thm:cholesky}), and close with the design effect ratio
classification and diagnostics. Throughout, we use the score functions
and Fisher information results from~\cref{app:identification}
(Supplementary Material~B) and the notation established in
\cref{sec:model}.

%% --- C.1: Pseudo-Likelihood and Separability ---
%% =============================================================================
%% sm_c1_pseudo.tex — SM-C, Subsection C.1: Pseudo-Likelihood and Separability
%% Body content only (\subsection level); parent structure in sm_c.tex
%% Primary source: dev/log/06_research-notes-survey-design-pseudo-posterior.qmd
%%   lines 570--660 (normalization, separability, weighted scores)
%% =============================================================================

\subsection{Pseudo-likelihood and separability}
\label{smc:pseudo-lik}

We summarize the pseudo-likelihood structure
from~\cref{sec:survey}.
By~\cref{prop:weighted-sep}, the decomposition
$\ell^{(w)}
 = \ell_{\mathrm{ext}}^{(w)}
 \allowbreak+ \ell_{\mathrm{int}}^{(w)}$
holds with weights
$\tilde{w}_i = w_i N / \sum_j w_j$.
We record three clarifying remarks.

\begin{remark}[Weight normalization]\label{smc:rem-normalization}
  Two normalization conventions appear in the survey literature.
  The \emph{sum-to-$N$} convention, $\tilde{w}_i = w_i N / \sum_j w_j$
  so that $\sum_i \tilde{w}_i = N$, is adopted throughout and is standard
  in the Bayesian pseudo-posterior literature
  \citep{SavitskyToth2016,WilliamsSavitsky2021}.
  The alternative \emph{sum-to-$\hat{M}$} convention retains unnormalized
  weights $\tilde{w}_i = w_i$, which for the NSECE gives
  $\sum_i w_i \approx 119{,}593$.  Under this convention the
  pseudo-likelihood resembles roughly $119{,}593$ observations rather
  than $N = 6{,}809$, leading to extreme overconcentration of the
  pseudo-posterior and credible intervals that are far too narrow.
  The sum-to-$N$ normalization ensures that the pseudo-likelihood is
  comparable in scale to an ordinary likelihood based on $N$ observations.
\end{remark}

\begin{remark}[Weighted score functions]\label{smc:rem-scores}
  The observation-level score functions for the extensive and intensive
  margins, along with their Fisher information and sign verification,
  are developed in~\cref{smb:scores} (Supplementary Material~B).
  The weighted versions used here simply multiply each per-observation
  score by $\tilde{w}_i$; the separability of
  \cref{prop:weighted-sep} guarantees that the weighted
  extensive-margin scores depend only on $(\balpha, \bdelta_{1,\cdot})$
  and the weighted intensive-margin scores depend only on
  $(\bbeta, \bdelta_{2,\cdot}, \kappa)$.
\end{remark}

\begin{remark}[Significance of separability under weighting]
\label{smc:rem-separability-significance}
  The separability result in~\cref{prop:weighted-sep} is not
  guaranteed for all survey adjustments.  Approaches that modify the
  likelihood structure---such as multilevel weight smoothing
  \citep{RabeHeskethSkrondal2006} or conditional likelihood
  methods---can break the additive structure.  Observation-level
  exponentiation $[f_i]^{\tilde{w}_i}$ preserves it precisely because
  $\tilde{w}_i$ is a scalar multiplier in the log scale and does
  not couple the two parameter blocks.  This separability is exploited
  in the sandwich estimator (\cref{sec:survey}), where the
  block-diagonality of $\mathbf{H}_{\mathrm{obs}}$ follows directly
  from the additive decomposition.
\end{remark}


%% --- C.2: Pseudo-Posterior Propriety ---
%% =============================================================================
%% sm_c2_propriety.tex — SM-C, Subsection C.2: Pseudo-Posterior Propriety
%% Body content only (\subsection level); parent structure in sm_c.tex
%% Primary source: dev/log/06_research-notes-survey-design-pseudo-posterior.qmd
%%   lines 741--781 (Definition 3.9, Theorem 3.1, Remarks 3.8--3.9)
%% =============================================================================

\subsection{Pseudo-posterior propriety}
\label{smc:propriety}

We now provide the complete proof of~\cref{thm:propriety}, which was
stated with a sketch in~\cref{sec:survey}.

\begin{proof}[Proof of \textup{\cref{thm:propriety}}]
  We verify that the pseudo-posterior
  \[
    \pi^{(w)}(\btheta \mid \by) \propto
    \exp\bigl[\ell^{(w)}(\btheta)\bigr]\,\pi(\btheta)
  \]
  is integrable over the parameter space.

  \textup{\textbf{Step 1} (Boundedness of the pseudo-likelihood).}
  For each observation $i$, the HBB probability satisfies
  $f_{\HBB}(y_i \mid \btheta) \in (0, 1]$ for all $\btheta$ in the
  interior of the parameter space.  Since $\tilde{w}_i > 0$, we have
  $f_i^{\tilde{w}_i} \le 1$ whenever $f_i \le 1$.  Therefore
  \[
    \exp\bigl[\ell^{(w)}(\btheta)\bigr]
    = \prod_{i=1}^{N} f_i^{\tilde{w}_i}
    \le 1
    \qquad \text{for all } \btheta.
  \]

  \textup{\textbf{Step 2} (Properness of the prior).}
  Each component of the prior $\pi(\btheta)$ is a proper distribution:
  \begin{itemize}[nosep]
    \item $\balpha, \bbeta \sim \Norm(\mathbf{0},\; 2^2\,\mathbf{I}_P)$
      \citep{GelmanEtAl2008}: proper with finite variance;
    \item $\boldsymbol{\tau} \sim \HalfNormal(0, 1)$ componentwise:
      proper on $\R_{>0}$;
    \item $\mathbf{R}_\varepsilon \sim \LKJ(2)$
      \citep{LewandowskiEtAl2009}: proper over the space of
      $2q \times 2q$ correlation matrices;
    \item $\log\kappa \sim \Norm(2,\; 1.5^2)$: proper.
  \end{itemize}
  The joint prior is a product of proper distributions, hence proper:
  $\int \pi(\btheta)\,d\btheta = 1$.

  \textup{\textbf{Step 3} (Integrability).}
  Combining Steps~1 and~2:
  \[
    \int \exp\bigl[\ell^{(w)}(\btheta)\bigr]\,\pi(\btheta)\,d\btheta
    \;\le\;
    \int 1 \cdot \pi(\btheta)\,d\btheta
    = 1 < \infty.
  \]
  Moreover, $\exp[\ell^{(w)}(\btheta)] > 0$ on the interior of the
  parameter space (since all probability terms $f_i$ are strictly
  positive there), so the pseudo-posterior is non-degenerate.
\end{proof}

\begin{remark}[Weight-invariance of propriety]
\label{smc:rem-weight-invariance}
  The proof above does not depend on the specific values or
  normalization of the weights---only on their positivity and
  finiteness.  Propriety therefore holds regardless of whether the
  weights are normalized to sum to~$N$, to sum to the population size
  $\hat{M}$, or under any other positive scaling.  However, the
  \emph{shape} and \emph{concentration} of the pseudo-posterior depend
  critically on the normalization convention, as discussed in
  \cref{smc:rem-normalization}: the sum-to-$\hat{M}$ convention
  produces extreme overconcentration, while the sum-to-$N$ convention
  preserves the information content of the actual sample.
\end{remark}

\begin{remark}[Contrast with improper priors]
\label{smc:rem-improper-priors}
  With improper (flat) priors on the regression coefficients,
  propriety would require the pseudo-likelihood to be integrable on
  its own---that is, $\int \exp[\ell^{(w)}(\btheta)]\,d\btheta < \infty$.
  In this setting the effective sample size (not the nominal sample
  size) governs integrability.  With $\ESS \approx 1{,}797$ and
  $\dim(\btheta) \le 606$, integrability would likely hold, but the
  formal verification is more delicate: it requires confirming that the
  weighted Hessian has full rank at the pseudo-MLE for every parameter
  block.  The use of proper priors sidesteps this issue entirely,
  ensuring propriety for \emph{any} positive weight configuration
  without conditions on the data.
\end{remark}


%% --- C.3: Bernstein--von Mises Approximation ---
%% =============================================================================
%% sm_c3_bvm.tex — SM-C, Subsection C.3: Bernstein–von Mises Approximation
%% Body content only (\subsection level); parent structure in sm_c.tex
%% Primary source: dev/log/06_research-notes-survey-design-pseudo-posterior.qmd
%%   lines 784--922 (Theorem 3.2, Proposition 3.8, Corollary 3.1, Remarks)
%% Correction applied: T1-1 (BvM reframe — diagnostic, not formal violation)
%% =============================================================================

\subsection{Bernstein--von Mises approximation}
\label{smc:bvm}

The sandwich correction in~\cref{sec:survey} rests on a
Bernstein--von Mises (BvM) theorem for pseudo-posteriors: as
$N \to \infty$, the pseudo-posterior converges to a normal distribution
centered at the pseudo-MLE with variance given by the inverse observed
Hessian, $\mathbf{H}_{\mathrm{obs}}^{-1}$.  We state the precise result, verify its regularity
conditions for the HBB model, and provide a finite-sample
diagnostic for the NSECE weight configuration that motivates the
sandwich correction as a robustness measure.

%% --- BvM Theorem ---

\begin{theorem}[Design-consistent BvM for the pseudo-posterior]
\label{smc:thm-bvm}
  Suppose the following regularity conditions hold\textup{:}
  \begin{enumerate}[label=\textup{(C\arabic*)},nosep]
    \item \label{smc:C1}%
      \textup{Smoothness.}
      The log-density $\log f_{\HBB}(y \mid \btheta)$ is twice
      continuously differentiable in $\btheta$ on the interior of the
      parameter space $\Theta$.
    \item \label{smc:C2}%
      \textup{Interior parameter.}
      The true superpopulation parameter $\btheta_0$ lies in the
      interior of $\Theta$.
    \item \label{smc:C3}%
      \textup{Weight regularity.}
      The ratio of the largest to smallest normalized weight satisfies
      $\max_i \tilde{w}_i / \min_i \tilde{w}_i = O(N^{\delta})$ for
      some $\delta < 1/2$ as $N \to \infty$.
    \item \label{smc:C4}%
      \textup{Information regularity.}
      The weighted Hessian $\mathbf{H}_{\mathrm{obs}}(\btheta_0)$ is
      positive definite, with eigenvalues bounded away from zero and
      infinity.
    \item \label{smc:C5}%
      \textup{Prior regularity.}
      The prior density $\pi(\btheta)$ is continuous and strictly
      positive in a neighborhood of $\btheta_0$.
  \end{enumerate}
  Then, as $N \to \infty$\textup{,}
  \begin{equation}\label{smc:eq-bvm}
    \sup_{B \in \cl{B}}
    \bigl|
      \pi^{(w)}(\btheta \in B \mid \by)
      - \Phi_{\hat{\btheta}^{(w)},\,\mathbf{H}_{\mathrm{obs}}^{-1}}(B)
    \bigr|
    \xrightarrow{P_{\btheta_0}} 0,
  \end{equation}
  where $\cl{B}$ is the class of Borel-measurable convex sets,
  $\hat{\btheta}^{(w)}$ is the pseudo-MLE,
  $\Phi_{\bmu,\bSigma}$ denotes the CDF of
  $\Norm(\bmu, \bSigma)$, and $\mathbf{H}_{\mathrm{obs}}$ is the
  observed Hessian of the pseudo-log-likelihood evaluated at
  $\hat{\btheta}^{(w)}$.
\end{theorem}

The result is a pseudo-posterior analogue of the classical BvM theorem
\citep{KleijnVanDerVaart2012}, extended to the survey-weighted setting by
\citet{SavitskyToth2016} and~\citet{WilliamsSavitsky2021}.  The key
difference from the standard BvM is condition~\ref{smc:C3}, which
constrains how rapidly the weight ratio may grow with sample size.
Importantly, the \emph{pseudo-posterior variance}
$\mathbf{H}_{\mathrm{obs}}^{-1}$ differs from the
\emph{design-consistent sampling variance} of $\hat{\btheta}^{(w)}$,
\[
  \hat{\btheta}^{(w)}
  \;\dot{\sim}\;
  \Norm\!\bigl(\btheta_0,\;
  \mathbf{V}_{\mathrm{sand}}\bigr),
  \qquad
  \mathbf{V}_{\mathrm{sand}}
  = \mathbf{H}_{\mathrm{obs}}^{-1}\,
    \mathbf{J}_{\mathrm{cluster}}\,
    \mathbf{H}_{\mathrm{obs}}^{-1},
\]
as defined in~\eqref{eq:V-sand}. Under informative sampling or model
misspecification,
$\mathbf{J}_{\mathrm{cluster}} \neq \mathbf{H}_{\mathrm{obs}}$,
so naive credible intervals based on $\mathbf{H}_{\mathrm{obs}}^{-1}$
are miscalibrated as a measure of frequentist uncertainty.  This
mismatch is precisely what motivates the sandwich correction developed
in~\cref{smc:sandwich}: \cref{smc:cor-coverage} below quantifies
the resulting coverage distortion, and
\cref{thm:cholesky} provides the operational calibration.

%% --- Verification of conditions ---

\paragraph{Verification.}
We verify each condition for the HBB model and the NSECE data.

\medskip\noindent\textit{\ref{smc:C1}: Smoothness.}\enspace
The HBB log-density involves $\log\Gamma$ functions (through the
beta-binomial PMF), the logistic link
$\sigma(\eta) = 1/(1+e^{-\eta})$, and the zero probability $p_0$
(a finite sum of $\log\Gamma$ terms).  Each component is $C^{\infty}$
on the interior of the parameter space, so the composite
$\log f_{\HBB}$ is smooth in $\btheta$.  \checkmark

\medskip\noindent\textit{\ref{smc:C2}: Interior parameter.}\enspace
The parameter space is
$\Theta = \R^P \times \R^P \times \R^{2QR} \times
\mathbb{S}_{++}^{2Q} \times \R_{>0}$,
and any finite parameter vector with $\kappa > 0$ and
$\bSigma_\varepsilon \succ 0$ lies in the interior.  \checkmark

\medskip\noindent\textit{\ref{smc:C3}: Weight regularity.}\enspace
This condition requires a dedicated analysis; see
\cref{smc:prop-weight-diagnostic} below.

\medskip\noindent\textit{\ref{smc:C4}: Information regularity.}\enspace
By~\cref{smb:prop-block-diag}, the Fisher information is
block-diagonal between the extensive and intensive margins.  For the
extensive block, positive definiteness follows from $\operatorname{rank}(\mathbf{X}) = P$
and $q_i \in (0,1)$ for all $i$ (standard logistic regression theory).
For the intensive block, separation of $\mu$ from $\kappa$ requires
variation in the denominators $n_i$
(\cref{smb:prop-mu-kappa-id}), which is satisfied empirically
($n_i$ ranges from~1 to~378 among positive responders).  \checkmark

\medskip\noindent\textit{\ref{smc:C5}: Prior regularity.}\enspace
The priors specified in~\cref{sec:priors} are all continuous and
strictly positive in the interior:
$\Norm(\mathbf{0},\, 2^2\,\mathbf{I}_P)$ on regression coefficients~\citep{GelmanEtAl2008},
$\HalfNormal(0,1)$ on scale parameters,
$\LKJ(2)$ on the correlation matrix~\citep{LewandowskiEtAl2009}, and
$\Norm(2,\, 1.5^2)$ on $\log\kappa$.  \checkmark

%% --- Finite-sample weight diagnostic ---

\begin{proposition}[Finite-sample weight diagnostic]
\label{smc:prop-weight-diagnostic}
  For the NSECE 2019 data with $N = 6{,}785$ and
  $w_{\max}/w_{\min} = 462$\textup{,} the weight ratio exceeds
  $N^{1/2} \approx 82.4$\textup{,} indicating that the sufficient
  condition \ref{smc:C3} is not comfortably satisfied for this design.
  The computation
  \begin{equation}\label{smc:eq-delta}
    \delta
    = \frac{\log 462}{\log 6{,}785}
    = \frac{6.136}{8.822}
    \approx 0.695
    > 0.5
  \end{equation}
  is a finite-sample diagnostic, not a valid asymptotic big-$O$ test
  of condition \ref{smc:C3}.  This motivates the sandwich variance
  correction as a robustness measure.
\end{proposition}

\begin{remark}[Interpretation of the diagnostic]
\label{smc:rem-weight-interpretation}
  Four observations clarify the scope of
  \cref{smc:prop-weight-diagnostic}.
  \begin{enumerate}[label=\textup{(\alph*)},nosep]
    \item Condition~\ref{smc:C3} is \emph{sufficient}, not necessary.
      The BvM approximation may remain adequate even when this bound is
      not met, provided the high-weight observations do not dominate
      the pseudo-likelihood.
    \item The fraction of observations with extreme weights is small:
      the top $1\%$ (weights $w_i > 100$) comprises fewer than 70
      observations out of $6{,}785$.  The diagnostic is driven by a
      small number of outlying weights, not a systematic pattern.
    \item Weight trimming can restore~\ref{smc:C3}.  Capping at the
      99th percentile reduces $w_{\max}$ to approximately~100, giving
      $\delta \approx 0.522$; capping at the 95th percentile
      ($w_{\max} \approx 50$) yields $\delta \approx 0.443 < 0.5$,
      restoring \ref{smc:C3}.  Weight trimming is examined as a
      sensitivity analysis in~\cref{app:extended-results}.
    \item In practice, the sandwich correction
      \eqref{eq:V-sand} provides the appropriate variance adjustment
      regardless of whether~\ref{smc:C3} holds exactly: it
      automatically accounts for the inflated variance due to extreme
      weights.
  \end{enumerate}
\end{remark}

%% --- Coverage calibration ---

\begin{corollary}[Coverage miscalibration of naive credible intervals]
\label{smc:cor-coverage}
  Under the conditions of~\cref{smc:thm-bvm}, define\textup{:}
  \begin{itemize}[nosep]
    \item $\mathrm{CI}_{\mathrm{naive}}(\alpha)$\textup{:} the
      $\alpha$-level credible interval from the pseudo-posterior
      marginal, with variance
      $[\mathbf{H}_{\mathrm{obs}}]^{-1}$\textup{;}
    \item $\mathrm{CI}_{\mathrm{sand}}(\alpha)
      = \hat{\theta}^{(w)} \pm z_{\alpha/2}\sqrt{V_{\mathrm{sand}}}$\textup{:}
      the sandwich-corrected interval.
  \end{itemize}
  Then\textup{:}
  \begin{enumerate}[label=\textup{(\alph*)},nosep]
    \item The naive interval has asymptotic coverage
      $\Phi(z_{\alpha/2}\, r) - \Phi(-z_{\alpha/2}\, r)$, where
      $r = \sqrt{\lambda_{\mathbf{H}} / \lambda_{\mathrm{sand}}}$ is
      the ratio of the naive to the sandwich standard deviation and
      $\Phi$ denotes the standard normal CDF.
    \item When the sampling is non-informative and the model is
      correctly specified, $\mathbf{J}_{\mathrm{cluster}} \approx
      \mathbf{H}_{\mathrm{obs}}$, so $r \approx 1$ and the naive
      interval has approximately correct coverage.
    \item Under informative sampling or model misspecification,
      $r \neq 1$.  If $r < 1$ \textup{(}naive variance too
      small\textup{)}, the interval is anti-conservative
      \textup{(}undercoverage\textup{)}---the typical case when the
      pseudo-posterior overconcentrates.  If $r > 1$, the interval is
      conservative.
  \end{enumerate}
\end{corollary}

\begin{proof}
  The BvM theorem (\cref{smc:thm-bvm}) implies that the
  pseudo-posterior converges to
  $\Norm(\hat{\btheta}^{(w)},\, \mathbf{H}_{\mathrm{obs}}^{-1})$,
  while the correct sampling distribution of $\hat{\btheta}^{(w)}$ is
  approximately $\Norm(\btheta_0,\, \mathbf{V}_{\mathrm{sand}})$.
  The ratio of marginal standard deviations determines the coverage
  distortion.
\end{proof}

\begin{remark}[Calibration via the Cholesky transformation]
\label{smc:rem-cholesky-calibration}
  \cref{thm:cholesky} provides the operational calibration
  strategy: the affine transformation
  \eqref{eq:cholesky-transform} rescales each MCMC draw so that
  the adjusted posterior has covariance
  $\mathbf{V}_{\mathrm{sand}}$ while preserving the posterior mean.
  This replaces the naive variance $\mathbf{H}_{\mathrm{obs}}^{-1}$
  with the design-consistent sandwich variance, correcting the
  coverage distortion quantified in~\cref{smc:cor-coverage}.  The
  block-wise implementation---applying the transformation separately
  to fixed effects, random effects, and hyperparameters---is discussed
  in~\cref{sec:survey}; the full proof appears in
  \cref{smc:cholesky} below.
\end{remark}


%% --- C.4: Sandwich Variance Estimator ---
%% =============================================================================
%% sm_c4_sandwich.tex — SM-C, Subsection C.4: Sandwich Variance Estimator
%% Body content only (\subsection level); parent structure in sm_c.tex
%% Primary source: dev/log/06_research-notes-survey-design-pseudo-posterior.qmd
%%   lines 1003--1195 (bread/meat matrices, cluster-robust J, SVC absorption)
%% Corrections applied: T2-3 ("exactly" -> "in expectation"), via ref to smb:rem-JeqH
%% =============================================================================

\subsection{Sandwich variance estimator}
\label{smc:sandwich}

The main text (\cref{sec:survey}) defines the bread matrix
$\mathbf{H}_{\mathrm{obs}}$~\eqref{eq:H-obs}, the cluster-robust meat
matrix $\mathbf{J}_{\mathrm{cluster}}$~\eqref{eq:J-cluster}, and the
sandwich variance $\mathbf{V}_{\mathrm{sand}}$~\eqref{eq:V-sand}.  This
subsection provides the detailed derivation: the block-diagonal
structure of~$\mathbf{H}_{\mathrm{obs}}$, the non-block-diagonal
structure of~$\mathbf{J}_{\mathrm{cluster}}$, the cluster-robust
aggregation formula with stratification~\citep{Binder1983}, and the partial absorption of
within-state clustering by the state-varying coefficients.  Throughout,
we use the score functions from~\cref{smb:scores} and the sandwich
convention defined below.

%% ---------------------------------------------------------------------------
%%  C.4.1  Bread matrix H
%% ---------------------------------------------------------------------------

\paragraph{Bread matrix.}
The bread matrix is the negative weighted observed Hessian at the
posterior mean $\hat{\btheta}$,
\begin{equation}\label{smc:eq-H}
  \mathbf{H}_{\mathrm{obs}}
  = -\sum_{i=1}^{N}\tilde{w}_i\,
    \frac{\partial^2\log f_{\HBB}(y_i \mid \btheta)}
         {\partial\btheta\,\partial\btheta\t}
    \bigg|_{\btheta = \hat{\btheta}}.
\end{equation}
By the separability of the hurdle log-likelihood
(\cref{smc:rem-separability-significance}) and the block-diagonal Fisher
information (\cref{smb:prop-block-diag}), the bread inherits a two-block
structure:

\begin{proposition}[Block-diagonal bread matrix]
\label{smc:prop-H-blockdiag}
  The observed Hessian decomposes as
  \begin{equation}\label{smc:eq-H-block}
    \mathbf{H}_{\mathrm{obs}}
    = \begin{pmatrix}
        \mathbf{H}_{\textup{ext}} & \mathbf{0} \\
        \mathbf{0} & \mathbf{H}_{\textup{int}}
      \end{pmatrix},
  \end{equation}
  where\textup{:}
  \begin{enumerate}[label=\textup{(\roman*)},nosep]
    \item The extensive block is the weighted logistic information
    \begin{equation}\label{smc:eq-H-ext}
      \mathbf{H}_{\textup{ext}}
      = \sum_{i=1}^{N}\tilde{w}_i\,q_i(1-q_i)\,
        \mathbf{d}_{1,i}\,\mathbf{d}_{1,i}\t,
    \end{equation}
    with $\mathbf{d}_{1,i}
    = \partial\eta_i^{(1)}/\partial(\balpha\t,\bdelta_{1,s[i]}\t)\t$.
    This matrix is non-stochastic: it depends on $q_i = q_i(\hat{\btheta})$
    but not on the observed data $z_i$.

    \item The intensive block is the weighted negative Hessian of the
    truncated beta-binomial log-likelihood,
    \begin{equation}\label{smc:eq-H-int}
      \mathbf{H}_{\textup{int}}
      = -\sum_{i \in \cl{I}_1}\tilde{w}_i\,
        \frac{\partial^2 \ell_i^{(2+)}}
             {\partial\boldsymbol{\phi}_2\,\partial\boldsymbol{\phi}_2\t}
        \bigg|_{\hat{\btheta}},
    \end{equation}
    where $\boldsymbol{\phi}_2
    = (\bbeta\t,\log\kappa,\bdelta_{2,s[i]}\t)\t$,
    $\cl{I}_1 = \{i : y_i > 0\}$, and
    $\ell_i^{(2+)} = \log f_{\textup{BB}}(y_i \mid n_i,\mu_i,\kappa)
     - \log(1-p_{0,i})$.
    This block depends on the observed counts $y_i$.
  \end{enumerate}
\end{proposition}

\begin{proof}
  By~\cref{smb:prop-block-diag}, the cross-derivative
  $\partial^2\ell_i/\partial\boldsymbol{\phi}_1\,\partial\boldsymbol{\phi}_2\t
  = \mathbf{0}$
  for every observation~$i$.  Multiplying by $\tilde{w}_i > 0$ and
  summing preserves this zero block, giving~\eqref{smc:eq-H-block}.
  The extensive-margin form~\eqref{smc:eq-H-ext} follows from
  $-\partial^2\ell_i^{(1)}/\partial(\eta_i^{(1)})^2 = q_i(1-q_i)$
  (\cref{smb:prop-ext-paramspace}) and the chain rule.
\end{proof}

%% ---------------------------------------------------------------------------
%%  C.4.2  Meat matrix J — individual level and non-block-diagonality
%% ---------------------------------------------------------------------------

\paragraph{Meat matrix and cross-margin coupling.}
The individual-level meat matrix is
\begin{equation}\label{smc:eq-J-ind}
  \mathbf{J}_{\textup{ind}}
  = \sum_{i=1}^{N}\tilde{w}_i^2\,\mathbf{s}_i\,\mathbf{s}_i\t,
\end{equation}
where $\mathbf{s}_i
= \partial\log f_{\HBB}(y_i \mid \btheta)/\partial\btheta
  \big|_{\hat{\btheta}}$
is the observation-level score.

\begin{remark}[Non-block-diagonal meat]
\label{smc:rem-J-nonblockdiag}
  Unlike the bread, the meat matrix is \emph{not} block-diagonal between
  the extensive and intensive margins.  The cross-block is
  \begin{equation}\label{smc:eq-J-cross}
    \mathbf{J}^{\textup{ext,int}}
    = \sum_{i=1}^{N}\tilde{w}_i^2\,
      \mathbf{s}_i^{\textup{ext}}\,
      (\mathbf{s}_i^{\textup{int}})\t.
  \end{equation}
  For $y_i = 0$, the intensive score $\mathbf{s}_i^{\textup{int}}
  = \mathbf{0}$, so the cross-product vanishes.  But for $y_i > 0$,
  both $\mathbf{s}_i^{\textup{ext}} = (1-q_i)\mathbf{d}_{1,i}$ and
  $\mathbf{s}_i^{\textup{int}} \neq \mathbf{0}$, so the cross-block
  receives a nonzero contribution from every positive observation.  In
  consequence, the sandwich variance
  $\mathbf{V}_{\mathrm{sand}}
  = \mathbf{H}_{\mathrm{obs}}^{-1}\mathbf{J}\mathbf{H}_{\mathrm{obs}}^{-1}$
  is not block-diagonal even though $\mathbf{H}_{\mathrm{obs}}$ is.  This
  cross-margin coupling is a general feature of sandwich estimators in
  hurdle models: a single PSU contributes weighted scores from both
  margins simultaneously.
\end{remark}

%% ---------------------------------------------------------------------------
%%  C.4.3  E[J] = H under equal weights
%% ---------------------------------------------------------------------------

\begin{remark}[Reduction under equal weights]
\label{smc:rem-equal-weights}
  Under uniform weights $\tilde{w}_i = 1$ and correct model
  specification, the information identity gives
  $\E[\mathbf{J}_{\textup{ind}}] = \mathbf{H}_{\mathrm{obs}}$
  \citep[see][for the pseudo-posterior analogue]{WilliamsSavitsky2021}, so
  $\mathbf{V}_{\mathrm{sand}} \to \mathbf{H}_{\mathrm{obs}}^{-1}$ and
  the sandwich reduces to the model-based variance.  For the extensive
  margin, $\E[\mathbf{J}_{\textup{ext}}] = \mathbf{H}_{\textup{ext}}$
  holds exactly, because the logistic Hessian is non-random (see
  \cref{smb:rem-ext-info}; under unequal survey weights the identity
  breaks because $\mathbf{H}$ involves $\tilde{w}_i$ while
  $\E[\mathbf{J}]$ involves $\tilde{w}_i^2$).  For the intensive margin, the identity holds only
  in expectation and the finite-sample discrepancy between
  $\mathbf{J}_{\textup{int}}$ and $\mathbf{H}_{\textup{int}}$ is
  nonzero.
\end{remark}

%% ---------------------------------------------------------------------------
%%  C.4.4  Cluster-robust meat (Proposition 4.2)
%% ---------------------------------------------------------------------------

\paragraph{Cluster-robust aggregation.}
The individual-level $\mathbf{J}_{\textup{ind}}$ treats observations as
independent.  In the NSECE stratified cluster design
\citep{NSECEProjectTeam2022}, providers within
the same primary sampling unit (PSU) share neighborhood
characteristics, inducing within-PSU correlation.  The cluster-robust
meat matrix replaces individual scores with PSU-level aggregates,
centered within strata.

\begin{definition}[PSU-level score totals]
\label{smc:def-psu-scores}
  For stratum $h = 1,\ldots,H$ and PSU $c = 1,\ldots,C_h$ within
  stratum~$h$, define
  \begin{equation}\label{smc:eq-s-psu}
    \bar{\mathbf{s}}_{hc}
    = \sum_{i \in \textup{PSU}(h,c)} \tilde{w}_i\,\mathbf{s}_i,
  \end{equation}
  where $\textup{PSU}(h,c) = \{i : \texttt{vstratum}[i]=h,\;
  \texttt{vpsu}[i]=c\}$.  The stratum mean is
  \begin{equation}\label{smc:eq-s-stratum}
    \bar{\mathbf{s}}_h
    = \frac{1}{C_h}\sum_{c=1}^{C_h}\bar{\mathbf{s}}_{hc}.
  \end{equation}
\end{definition}

\begin{proposition}[Cluster-robust meat matrix]
\label{smc:prop-cluster-robust}
  Under the stratified cluster design with $H$~strata and
  $C_h$~PSUs in stratum~$h$, the cluster-robust meat matrix is
  \begin{equation}\label{smc:eq-J-cluster}
    \boxed{%
      \mathbf{J}_{\mathrm{cluster}}
      = \sum_{h=1}^{H}\frac{C_h}{C_h-1}
        \sum_{c=1}^{C_h}
        (\bar{\mathbf{s}}_{hc} - \bar{\mathbf{s}}_h)\,
        (\bar{\mathbf{s}}_{hc} - \bar{\mathbf{s}}_h)\t.
    }
  \end{equation}
  This estimator accounts for\textup{:}
  \textup{(i)}~within-PSU correlation, by aggregating scores to PSU
  totals\textup{;}
  \textup{(ii)}~stratification, by centering PSU scores within strata
  so that only within-stratum variation across PSUs contributes\textup{;}
  \textup{(iii)}~unequal cluster sizes, through the weight-aggregated
  PSU totals $\bar{\mathbf{s}}_{hc}$\textup{;}
  \textup{(iv)}~finite-sample degrees of freedom, via the factor
  $C_h/(C_h-1)$.
\end{proposition}

\begin{proof}
  The design-based variance of the weighted score total under
  stratified sampling decomposes as
  \[
    \Var_{\textup{design}}\!\Bigl(\sum_i\tilde{w}_i\,\mathbf{s}_i\Bigr)
    = \sum_{h=1}^{H}
      \Var_h\!\Bigl(\sum_{c=1}^{C_h}\bar{\mathbf{s}}_{hc}\Bigr),
  \]
  because sampling across strata is independent.  Within stratum~$h$,
  PSUs are sampled without replacement from the population of
  $C_h^{\ast}$~PSUs.  Under the ``ultimate cluster'' assumption that
  $C_h / C_h^{\ast} \approx 0$ (negligible PSU sampling fraction),
  the finite-population correction is dropped and the within-stratum
  variance estimator is
  \[
    \widehat{\Var}_h
    = \frac{C_h}{C_h-1}
      \sum_{c=1}^{C_h}
      (\bar{\mathbf{s}}_{hc}-\bar{\mathbf{s}}_h)\,
      (\bar{\mathbf{s}}_{hc}-\bar{\mathbf{s}}_h)\t,
  \]
  the standard unbiased estimator of the variance of a total from
  $C_h$~units.  Summing over strata gives
  $\mathbf{J}_{\mathrm{cluster}}$.
\end{proof}

\begin{remark}[Degrees of freedom]
\label{smc:rem-dof}
  The cluster-robust estimator has
  $\sum_{h=1}^{H}(C_h - 1) = 415 - 30 = 385$ effective degrees of
  freedom.  With $C_h$ ranging from~5 to~36 (median~$\approx 14$), each
  stratum contributes at least~4 degrees of freedom.  This comfortably
  exceeds the rule-of-thumb threshold of $\sum_h(C_h - 1) > 30$ for
  stable sandwich variance estimation.
\end{remark}

\begin{remark}[Individual versus cluster-level $\mathbf{J}$]
\label{smc:rem-J-comparison}
  When there is no within-PSU correlation conditional on covariates,
  $\mathbf{J}_{\mathrm{cluster}}$ and $\mathbf{J}_{\textup{ind}}$
  converge to the same population quantity.  In practice, within-PSU
  correlation inflates $\mathbf{J}_{\mathrm{cluster}}$ relative to
  $\mathbf{J}_{\textup{ind}}$; the parameter-specific ratio
  $\diag(\mathbf{V}_{\mathrm{sand}})
  / \diag(\mathbf{H}_{\mathrm{obs}}^{-1}
           \mathbf{J}_{\textup{ind}}
           \mathbf{H}_{\mathrm{obs}}^{-1})$
  provides an estimate of the cluster design effect.
  The design effect ratio (DER) classification and diagnostics are
  developed in~\cref{smc:der} below.
\end{remark}

%% ---------------------------------------------------------------------------
%%  C.4.5  SVC partial absorption (Proposition 4.3)
%% ---------------------------------------------------------------------------

\begin{proposition}[Partial absorption of clustering by SVCs]
\label{smc:prop-svc-absorption}
  The state-varying coefficients $\bepsilon_s$ partially absorb
  within-state clustering.  For any two providers $i,j$ in the same
  state~$s$ and the same PSU,
  \begin{equation}\label{smc:eq-svc-absorption}
    \E_{\bepsilon_s}\!\bigl[
      \Cov(s_i,\,s_j \mid \bepsilon_s)
    \bigr]
    \;\leq\;
    \Cov(s_i,\,s_j),
  \end{equation}
  with equality if and only if $\Cov(\E[s_i \mid \bepsilon_s],\,
  \E[s_j \mid \bepsilon_s]) = 0$.
\end{proposition}

\begin{proof}
  By the law of total covariance,
  \[
    \Cov(s_i,\,s_j)
    = \E\!\bigl[\Cov(s_i,\,s_j \mid \bepsilon_s)\bigr]
    + \Cov\!\bigl(
        \E[s_i \mid \bepsilon_s],\;
        \E[s_j \mid \bepsilon_s]
      \bigr).
  \]
  The second term is non-negative, so
  $\Cov(s_i,s_j) \geq \E[\Cov(s_i,s_j \mid \bepsilon_s)]$.
  At the posterior mean $\hat{\btheta}$ (which includes state-specific
  $\hat{\bepsilon}_s$), the plug-in conditional covariance is an estimate
  of $\Cov(s_i,s_j \mid \bepsilon_s)$; hence
  $\mathbf{J}_{\mathrm{cluster}}$ evaluated at $\hat{\btheta}$ will
  be smaller than under a model without SVCs.
\end{proof}

\begin{remark}[Residual clustering]
\label{smc:rem-residual-clustering}
  The inequality~\eqref{smc:eq-svc-absorption} shows that SVCs reduce
  but do not eliminate within-PSU correlation.  The residual correlation
  reflects neighborhood-level effects not captured by the five
  state-varying covariates (poverty, race/ethnicity, urbanicity): local
  labor market conditions, municipal zoning, co-located service
  providers, and other spatially structured unobservables.  The sandwich
  correction with cluster-robust~$\mathbf{J}_{\mathrm{cluster}}$ remains
  necessary to account for this residual design effect.
\end{remark}


%% --- C.5: Cholesky Affine Transformation ---
%% =============================================================================
%% sm_c5_cholesky.tex — SM-C, Subsection C.5: Cholesky Affine Transformation
%% Body content only (\subsection level); parent structure in sm_c.tex
%% Primary source: dev/log/06_research-notes-survey-design-pseudo-posterior.qmd
%%   lines 1197--1261 (theorem proof, block-wise application, limitations)
%% Corrections applied: T1-4 (MCMC vs design-consistent shrinkage),
%%   T2-1 (block-wise Cholesky caveat)
%% =============================================================================

\subsection{Cholesky affine transformation}
\label{smc:cholesky}

The sandwich variance $\mathbf{V}_{\mathrm{sand}}$~\eqref{eq:V-sand}
provides the design-consistent covariance, but the MCMC draws
$\{\btheta^{(m)}\}_{m=1}^M$ from the pseudo-posterior have empirical
covariance $\bSigma_{\mathrm{MCMC}} \approx \mathbf{H}_{\mathrm{obs}}^{-1}$,
not $\mathbf{V}_{\mathrm{sand}}$.  In a frequentist setting one would simply
report $\mathbf{V}_{\mathrm{sand}}$ as the variance estimator; in the
Bayesian framework, we require the posterior draws themselves to have the
correct covariance so that credible intervals are properly calibrated
\citep{WilliamsSavitsky2021}.
\cref{thm:cholesky} in the main text states the affine transformation
that achieves this.  We now provide the full proof.

%% ---------------------------------------------------------------------------
%%  Proof of Theorem (thm:cholesky)
%% ---------------------------------------------------------------------------

\begin{proof}[Proof of \textup{\cref{thm:cholesky}}]
  Define centered draws
  $\mathbf{u}^{(m)} = \btheta^{(m)} - \hat{\btheta}$, so that
  $\bSigma_{\mathrm{MCMC}}
  = M^{-1}\sum_{m=1}^{M}\mathbf{u}^{(m)}(\mathbf{u}^{(m)})\t$.
  Set
  \begin{equation}\label{smc:eq-A}
    \mathbf{A}
    = \mathbf{L}_{\mathrm{sand}}\,\mathbf{L}_{\mathrm{MCMC}}^{-1},
  \end{equation}
  where $\mathbf{L}_{\mathrm{MCMC}}
  = \mathrm{chol}(\bSigma_{\mathrm{MCMC}})$ and
  $\mathbf{L}_{\mathrm{sand}}
  = \mathrm{chol}(\mathbf{V}_{\mathrm{sand}})$ are lower-triangular
  Cholesky factors.  The transformation~\eqref{eq:cholesky-transform} is
  then $\btheta^{*(m)} = \hat{\btheta} + \mathbf{A}\,\mathbf{u}^{(m)}$.

  \medskip\noindent
  \textup{(a) Mean preservation.}
  \[
    \frac{1}{M}\sum_{m=1}^{M}\btheta^{*(m)}
    = \hat{\btheta}
      + \mathbf{A}\,\Bigl(\frac{1}{M}\sum_{m=1}^{M}\mathbf{u}^{(m)}\Bigr)
    = \hat{\btheta} + \mathbf{A}\cdot\mathbf{0}
    = \hat{\btheta}.
  \]

  \noindent
  \textup{(b) Covariance recovery.}
  \begin{align}
    \frac{1}{M}\sum_{m=1}^{M}
      (\btheta^{*(m)} - \hat{\btheta})\,
      (\btheta^{*(m)} - \hat{\btheta})\t
    &= \mathbf{A}\,\bSigma_{\mathrm{MCMC}}\,\mathbf{A}\t
      \notag \\
    &= \mathbf{L}_{\mathrm{sand}}\,
       \mathbf{L}_{\mathrm{MCMC}}^{-1}\,
       \mathbf{L}_{\mathrm{MCMC}}\,
       \mathbf{L}_{\mathrm{MCMC}}\t\,
       \mathbf{L}_{\mathrm{MCMC}}^{-\t}\,
       \mathbf{L}_{\mathrm{sand}}\t
      \notag \\
    &= \mathbf{L}_{\mathrm{sand}}\,\mathbf{I}\,
       \mathbf{L}_{\mathrm{sand}}\t
     = \mathbf{V}_{\mathrm{sand}}.
      \label{smc:eq-cholesky-var}
  \end{align}

  \noindent
  \textup{(c) Affine equivariance.}
  The standardized draws
  $\mathbf{z}^{(m)} = \mathbf{L}_{\mathrm{MCMC}}^{-1}\mathbf{u}^{(m)}$
  satisfy
  \[
    M^{-1}\sum_m \mathbf{z}^{(m)}(\mathbf{z}^{(m)})\t = \mathbf{I}.
  \]
  The transformation maps
  \[
    \btheta^{*(m)} - \hat{\btheta}
    = \mathbf{L}_{\mathrm{sand}}\,\mathbf{z}^{(m)},
  \]
  which is a linear rescaling of the standardized draws.  The
  Wald confidence intervals---computed as
  $\hat{\theta}_p \pm z_{\alpha/2}\sqrt{[\mathbf{V}_{\mathrm{sand}}]_{pp}}$---depend
  only on the marginal variances of~$\btheta^*$ and are therefore
  invariant to the choice of~$\mathbf{A}$.  Marginal quantiles of
  the transformed draws may differ from those of the original
  draws when~$\mathbf{A}$ is non-diagonal.
\end{proof}

%% ---------------------------------------------------------------------------
%%  Block-wise application
%% ---------------------------------------------------------------------------

\begin{remark}[Block-wise application]
\label{smc:rem-blockwise}
  In practice the transformation~\eqref{eq:cholesky-transform} is applied
  block-wise rather than to the full
  $\approx\!616$-dimensional vector simultaneously.  The three
  parameter classes are listed below; the numerical Cholesky procedure
  operates on blocks~(i) and~(ii) only, since block~(iii) receives no
  correction
  (\cref{rem:blockwise} in the main text).  The rationale follows from
  \cref{smc:prop-H-blockdiag} and~\cref{smc:rem-J-nonblockdiag}\textup{:}
  \begin{enumerate}[label=\textup{(\roman*)},nosep]
    \item \emph{Fixed effects}
      $\boldsymbol{\xi}
      = (\balpha\t,\bbeta\t,\log\kappa)\t \in \R^{11}$\textup{:}
      full Cholesky correction.  The $\DER$ ranges from $1.14$ to
      $4.18$ (mean $2.11$), and the corresponding confidence-interval
      inflation factors are $\sqrt{\DER} \in [1.07, 2.04]$.
    \item \emph{State random effects}
      $\bepsilon_s \in \R^{10}$, $s = 1,\ldots,S$\textup{:} partial
      correction.  The posterior variance of $\bepsilon_s$ mixes the
      likelihood information from providers in state~$s$ (affected by
      survey weights) with the shrinkage toward the prior mean
      $\bGamma_k\mathbf{v}_s$ (unaffected).  When the prior
      dominates (small states), $\DER_s \approx 1$ and no correction
      is needed; when the likelihood dominates (large states),
      $\DER_s$ can exceed~$3$ and the correction is substantial.
    \item \emph{Hyperparameters}
      $(\bGamma_k, \bSigma_\varepsilon)$\textup{:} no correction.
      The sandwich estimator is not applicable to variance components
      estimated from between-state variation; the DER concept does not
      extend to these parameters.
  \end{enumerate}
\end{remark}

%% ---------------------------------------------------------------------------
%%  T2-1: Approximation quality of block-wise approach
%% ---------------------------------------------------------------------------

\begin{remark}[Approximation quality of the block-wise correction]
\label{smc:rem-blockwise-quality}
  The block-wise application ignores the cross-block covariance between
  fixed effects $\boldsymbol{\xi}$ and random effects $\bepsilon_s$ in
  $\mathbf{V}_{\mathrm{sand}}$.  The approximation is essentially exact
  when the prior dominates ($\ESS_s \ll q$), since $\bepsilon_s$ is then
  nearly independent of the data.  When data dominate ($\ESS_s \gg q$),
  the cross-block covariance can be non-negligible; the sensitivity
  comparison in~\cref{app:extended-results} addresses this case.
\end{remark}

%% ---------------------------------------------------------------------------
%%  T1-4: MCMC vs design-consistent shrinkage
%% ---------------------------------------------------------------------------

\begin{remark}[MCMC versus design-consistent shrinkage]
\label{smc:rem-mcmc-vs-design}
  The shrinkage factor
  $B_s^{(w)}
  = \bSigma_\varepsilon(\bSigma_\varepsilon + \mathbf{V}_s^{-1})^{-1}$
  describes the \emph{design-consistent} (post-sandwich-correction)
  posterior behavior, not the raw MCMC pseudo-posterior.  The MCMC
  sampler sees precision proportional to
  $W_s = \sum_{i \in s}\tilde{w}_i$, which exceeds the effective
  sample size $\ESS_s = W_s^2 / \sum_{i \in s}\tilde{w}_i^2$ whenever
  within-state weights are heterogeneous.  The Cholesky affine
  transformation~\eqref{eq:cholesky-transform} maps the MCMC output to
  the design-consistent target, inflating the variance by the factor
  $\DER_s$ (\cref{smc:der}).  The gap $W_s - \ESS_s$ quantifies the
  over-concentration of the raw pseudo-posterior relative to the
  design-aware target\textup{:} large gaps (common in states with a
  few heavily weighted providers) signal that the naive MCMC credible
  intervals are too narrow and that the sandwich correction is most
  important.
\end{remark}

%% ---------------------------------------------------------------------------
%%  Limitations
%% ---------------------------------------------------------------------------

\begin{remark}[Non-Gaussian posterior shapes]
\label{smc:rem-nongaussian}
  The Cholesky transformation is an affine operation that preserves
  only first and second moments.  If the pseudo-posterior is
  substantially non-Gaussian (skewed or multimodal), the transformed
  draws will have the correct covariance but may misrepresent the
  posterior shape.  For the HBB model, the fixed effects
  $(\balpha,\bbeta)$ are expected to be approximately Gaussian by the
  Bernstein--von Mises result (\cref{smc:thm-bvm}) given
  $N = 6{,}809$; the log-dispersion $\log\kappa$ is also
  well-approximated by a normal.  For variance components
  ($\bSigma_\varepsilon$), the posterior may exhibit right-skewness,
  but these parameters are not subject to the Cholesky correction
  (Block~3 above).
\end{remark}

\begin{remark}[Positive definiteness requirements]
\label{smc:rem-posdef}
  The transformation requires both $\bSigma_{\mathrm{MCMC}}$ and
  $\mathbf{V}_{\mathrm{sand}}$ to be positive definite so that their
  Cholesky factors exist.  For $\bSigma_{\mathrm{MCMC}}$, positive
  definiteness holds almost surely when the number of MCMC draws
  satisfies $M \gg \dim(\boldsymbol{\xi}) = 11$; in our application
  $M = 4{,}000$ post-warmup draws.  For
  $\mathbf{V}_{\mathrm{sand}}$, positive definiteness requires the
  effective degrees of freedom to exceed the parameter dimension:
  $\sum_{h=1}^{H}(C_h - 1) = 386 > 11$ (\cref{smc:rem-dof}), which
  is easily satisfied.
\end{remark}


%% --- C.6: Design Effect Ratio ---
%% =============================================================================
%% sm_c6_der.tex — SM-C, Subsection C.6: Design Effect Ratio
%% Body content only (\subsection level); parent structure in sm_c.tex
%% Primary source: dev/log/06_research-notes-survey-design-pseudo-posterior.qmd
%%   lines 1262--1291 (Definition 4.4, Proposition 4.4, Remark 4.12)
%% Corrections applied: T2-2 (DER heuristic label + overestimation),
%%   T2-4 (intercept identifiability)
%% =============================================================================

\subsection{Design effect ratio classification and diagnostics}
\label{smc:der}

The sandwich variance $\mathbf{V}_{\mathrm{sand}}$~\eqref{eq:V-sand}
and the Cholesky transformation~\eqref{eq:cholesky-transform} need not
be applied uniformly to every parameter.  The design effect ratio
(DER), defined in the main text~\eqref{eq:der}, provides a
parameter-specific diagnostic for the magnitude of the design effect:
\begin{equation}\label{smc:eq-der-recall}
  \DER_{p}
  = \frac{[\mathbf{V}_{\mathrm{sand}}]_{pp}}
         {[\mathbf{H}_{\mathrm{obs}}^{-1}]_{pp}}.
\end{equation}
When $\DER_p \approx 1$, the model-based variance already
approximates the design-consistent variance and no correction is
needed; when $\DER_p$ is substantially larger than~1, the Cholesky
correction is essential for calibrated inference.

%% ---------------------------------------------------------------------------
%%  DER classification
%% ---------------------------------------------------------------------------

\begin{proposition}[DER classification of parameters]
\label{smc:prop-der-classification}
  The parameters of the HBB model fall into three categories with
  respect to the Cholesky correction\textup{:}
  \begin{enumerate}[label=\textup{(\roman*)},nosep]
    \item \emph{Fixed effects}
      $\boldsymbol{\xi}
      = (\balpha\t,\bbeta\t,\log\kappa)\t$\textup{:}
      full correction required.  $\DER_p$ reflects the joint effect of
      unequal weights and within-PSU clustering on the pseudo-likelihood
      information.  Empirically, $\DER_p \in [1.14,\,4.18]$
      (mean~$2.11$), consistent with the Kish $\DEFF = 3.79$.

    \item \emph{State random effects}
      $\bepsilon_s$\textup{:} partial correction.
      The posterior variance of $\bepsilon_s$ blends
      likelihood-based information (affected by the design) with
      shrinkage toward the prior mean $\bGamma_k\bv_s$ (unaffected).
      States in which the prior dominates ($\ESS_s \ll q$) have
      $\DER_s \approx 1$; states in which the data dominate
      ($\ESS_s \gg q$) may have $\DER_s > 3$.

    \item \emph{Hyperparameters}
      $(\bGamma_k, \bSigma_\varepsilon)$\textup{:} no correction.
      The sandwich estimator is not applicable to variance components
      estimated from between-state variation; the DER concept does not
      extend to these parameters. Posterior intervals for
      hyperparameters are therefore model-based summaries whose
      calibration depends on correct specification of the hierarchical
      distribution.
  \end{enumerate}
\end{proposition}

\begin{proof}
  Part~(i) follows because the fixed-effect block of
  $\mathbf{H}_{\mathrm{obs}}$ and $\mathbf{J}_{\mathrm{cluster}}$
  depends directly on the observation-level weights $\tilde{w}_i$ and
  PSU aggregation, so $\DER_p$ is determined by the design effect.
  Part~(ii) follows from the posterior precision decomposition: for
  the normal--normal conjugate approximation, the posterior precision
  of $\bepsilon_s$ is the sum of the data precision (design-affected)
  and the prior precision (design-free).  When the prior precision
  dominates, the posterior variance is insensitive to the design, and
  $\DER_s \approx 1$.
  Part~(iii) holds because $\bGamma_k$ and $\bSigma_\varepsilon$ are
  identified from the between-state variation of the $S = 51$
  state-level random effects.  The survey design affects
  \emph{within-state} estimation but does not alter the between-state
  variance structure; hence the sandwich correction has no natural
  analogue for these parameters.
\end{proof}

%% ---------------------------------------------------------------------------
%%  DER heuristic for random effects (T2-2)
%% ---------------------------------------------------------------------------

\noindent
\textbf{Heuristic C.1} (Linear interpolation for random-effect DER).
\label{smc:heuristic-der-re}%
With shrinkage factor
$\lambda_{s,pp}
= \sigma_\varepsilon^2 / (\sigma_\varepsilon^2 + \sigma_e^2 / \ESS_s)$,
the DER for $\varepsilon_{s,p}$ is approximately
\begin{equation}\label{smc:eq-der-heuristic}
  \DER_{s,p}
  \;\approx\;
  1 + \lambda_{s,pp}\,(\DEFF_s - 1),
\end{equation}
interpolating between $\DER_{s,p} = 1$ (strong shrinkage) and
$\DER_{s,p} = \DEFF_s$ (no shrinkage).

\begin{remark}[Overestimation of the linear heuristic]
\label{smc:rem-heuristic-overestimate}
  {\upshape\textbf{(Correction T2-2.)}}
  The linear interpolation~\eqref{smc:eq-der-heuristic} is a
  \emph{heuristic}, not an exact result.  In the one-dimensional
  normal--normal conjugate model, the exact DER is
  \begin{equation}\label{smc:eq-der-exact}
    \DER_{s,p}^{\textup{exact}}
    = \frac{1 + \lambda_{s,pp}\,\DEFF_s}
           {(1 + \lambda_{s,pp})^2},
  \end{equation}
  which is always $\leq$~\eqref{smc:eq-der-heuristic}.  The
  discrepancy is largest at intermediate shrinkage
  ($\lambda_{s,pp} \approx 0.5$), where the linear heuristic can
  overestimate by up to~50\%.  For practical purposes, the heuristic
  remains useful as a conservative upper bound that quickly
  identifies parameters needing correction; the exact
  formula~\eqref{smc:eq-der-exact} can be computed alongside as a
  more accurate benchmark.
\end{remark}

%% ---------------------------------------------------------------------------
%%  DER diagnostic thresholds
%% ---------------------------------------------------------------------------

\begin{proposition}[DER diagnostic thresholds]
\label{smc:prop-der-diagnostic}
  As a post-estimation diagnostic, compute $\DER_p$ for each
  parameter and classify as follows\textup{:}
  \begin{itemize}[nosep]
    \item $\DER_p \in [0.8,\;1.2]$\textup{:}
      \emph{design-insensitive}.  The naive credible interval is
      adequately calibrated; sandwich correction is optional.
    \item $\DER_p \in (1.2,\;1.5]$\textup{:}
      \emph{moderate inflation}.  The confidence-interval width
      increases by $\sqrt{1.5} \approx 22\%$; correction is
      recommended.
    \item $\DER_p > 1.5$\textup{:}
      \emph{design-sensitive}.  The sandwich-corrected variance should
      be used as the primary inferential summary.
  \end{itemize}
  In the NSECE application, all 11 fixed-effect parameters have
  $\DER_p > 1.14$, placing them in the moderate-to-sensitive range.
\end{proposition}

%% ---------------------------------------------------------------------------
%%  Intercept identifiability remark (T2-4)
%% ---------------------------------------------------------------------------

\begin{remark}[Intercept identifiability under DER correction]
\label{smc:rem-intercept-id}
  {\upshape\textbf{(Correction T2-4.)}}
  Because the policy moderator vector $\bv_s$ includes an intercept
  column (see \cref{sec:policy}), the fixed
  intercept~$\alpha_1$ and the first column of $\bGamma_k$ have
  partially overlapping roles: the state-specific intercept is
  $\alpha_1 + \gamma_{k,1,1} + \varepsilon_{k,1,s}$.  The DER
  for~$\alpha_1$ includes variance attributable to both the design
  effect on the likelihood and the partial confounding with
  $\gamma_{k,1,1}$.  Formally, the decomposition
  \begin{equation}\label{smc:eq-intercept-decomp}
    [\mathbf{V}_{\mathrm{sand}}]_{\alpha_1,\alpha_1}
    = [\mathbf{V}_{\mathrm{sand}}^{\textup{design}}]_{\alpha_1,\alpha_1}
    + [\mathbf{V}^{\textup{id}}]_{\alpha_1,\alpha_1}
  \end{equation}
  is not cleanly separable; the DER should be interpreted with caution
  for intercept parameters.  The soft identification provided by the
  tighter prior $\Norm(0, 0.5^2)$ on the intercept column of
  $\bGamma_k$ (see \cref{sec:priors}) mitigates this issue by
  attributing the cross-state mean to~$\alpha_1$ and treating
  deviations as random effects.
\end{remark}

%% ---------------------------------------------------------------------------
%%  Computational savings
%% ---------------------------------------------------------------------------

\begin{remark}[Computational savings from DER classification]
\label{smc:rem-computational-savings}
  The three-way classification implies that the Cholesky
  transformation~\eqref{eq:cholesky-transform} need only be applied
  to the fixed-effect block ($\dim\boldsymbol{\xi} = 11$) and, for
  the largest states, the random-effect blocks.  This reduces the
  dimensionality of the Cholesky factorization from the full parameter
  dimension ($\approx\!616$) to at most $11$--$20$, yielding
  substantial computational savings.  The empirical DER values are
  reported in the main text (\cref{tab:fixed-effects}).
\end{remark}

